% Cover letter for the Knowledge-Based Systems submission.
% Compile: pdflatex cover-letter.tex   (no bibtex, no second pass needed)
\documentclass[10pt,a4paper]{article}
\usepackage[T1]{fontenc}
\usepackage[utf8]{inputenc}
\usepackage[margin=1.7cm]{geometry}
\usepackage{url}
\usepackage[hidelinks]{hyperref}
\pagestyle{empty}
\setlength{\parindent}{0pt}
\setlength{\parskip}{2.0pt}

\begin{document}

\begin{flushright}
Suxing Liu (corresponding author)\\
Department of IT Engineering, Mokwon University, Daejeon 35349, Korea\\
School of Digital Media Technology, Jiangxi Arts \& Ceramics Technology\\
Institute, Jingdezhen 333001, Jiangxi, China\\
\href{mailto:bentondoucet@gmail.com}{bentondoucet@gmail.com}
\end{flushright}

\vspace{4pt}
\today

\vspace{4pt}
The Editors-in-Chief\\
\emph{Knowledge-Based Systems}

\vspace{2pt}
\textbf{Re: submission of ``The Model Knows the Textbook but Not Your Class:
Curriculum Knowledge Graph Constraints for Slide Generation in Two
Vocational-Education Courses''}

\vspace{2pt}
Dear Editors,

On behalf of all authors, I am submitting the manuscript above for consideration
as a research article in \emph{Knowledge-Based Systems}.

\textbf{Where the paper sits in the journal's scope.} The journal's stated topics
include \emph{knowledge presentation and engineering}, and its objectives ask for
work that balances theory with practical study and develops knowledge-based
models, methods and tools with applications in education. This manuscript is a
knowledge-representation paper before it is a generation paper, and it makes two
claims of different kinds. An \emph{effect} claim: a graph-compiled contract
improves generated decks along named dimensions. A \emph{scope} claim: which
dimensions of a syllabus can be compiled into requirements that are simultaneously
generation constraints and acceptance criteria is a measurable property of the
representation rather than a matter of taste---and that boundary depends jointly
on the representation and on the backbone's willingness to execute a structural
instruction. The second is the claim we expect to outlive the course, and it is why
the generator is swapped for other models.

\textbf{The knowledge-engineering claim.} Course-level knowledge---the concepts a
session must cover, the misconception this cohort keeps committing, what must not
be introduced yet---is represented as a curriculum knowledge graph and compiled,
before generation, into an executable contract that then serves as the acceptance
criterion, so requirement and check are one data object and cannot drift. This
differs from retrieval: a retrieved passage decides what the model \emph{reads},
whereas a compiled contract decides what the output \emph{owes}, and an unmet
obligation is a defect by construction. Section~1.5 gives the representational
argument: presence, absence and a corrective slot admit hard constraints; order is
expressible in the graph but cannot be withheld by the interface.

\textbf{What the evidence supports.} The central comparison is not ``text present
vs absent'' but a $2\times2$ over the \emph{same} passage. Holding the
misconception statement fixed, the study crosses whether the contract carries a
slot for it: $1.000$ with structure and text, $0.111$ with structure alone,
$0.556$ with text alone, $0.222$ with neither. The interaction is $+0.556$---the
same passage is worth $+0.889$ inside the contract and $+0.334$ outside it---which
is the paper's claim in one line: the contract does not deliver the text, it makes
the text land. That is the isolation of a knowledge component the journal asks
for, and it is offered as a descriptive mechanism claim, not a significance claim:
with three sessions an exact sign-flip test cannot return a two-sided $p$ below
$0.25$, and the four cells are conditions run for other purposes rather than a
crossing randomised inside one batch.

Around it: a six-condition experiment in one higher-vocational course (54 runs),
three one-factor ablations (27 runs), a length-matched placebo, a blind review by
three teachers, the claim-bearing conditions replicated on two further models and
on a second course in another language, a run that holds the course fixed while
changing only the endpoint, and a bounded repair loop that takes the main grid from
6/9 to 9/9 in three calls while the weakest backbone moves only from 4/6 to 5/6.
The ablation attributes each remaining gain to a named block: the
forbidden-term list buys 0.083 of out-of-graph control and moves nothing else, and
the ordering row is reported as a null manipulation rather than as a price.
Seven measurement defects that had favoured
constraint-based methods are documented with their fixes; the paper states plainly
that it is a two-course case study with no learning-outcome measurement.

\textbf{Reproducibility.} Code, both curriculum graphs, prompts, the protocol and
all raw outputs are archived in a public repository (Zenodo DOI:
\url{https://doi.org/10.5281/zenodo.XXXXXXX}); every table here and in the
supplementary material is regenerated from it by the included scripts. No
proprietary or licensed textbook content is redistributed.

\textbf{Declarations.} The manuscript is original, is not under consideration
elsewhere, and all four authors have approved the version submitted and agree to
its submission here. No competing interests; no specific grant from funding agencies in the public,
commercial or not-for-profit sectors. The course data are aggregate classroom
statistics; no student-level record is included.

I would be glad to provide any further information the editors require.

\vspace{2pt}
Yours sincerely,

Suxing Liu\\
Corresponding author, on behalf of all authors

\end{document}
